\section{Continuity of the parent-Hamiltonian gap}
\label{sec:parent_hamiltonian_gap_continuity}

Along a continuous path of tensors, the phase argument of
\cite[Appendix~A, lines~2475--2580]{Schuch2011Phases} needs a gap of the
parent Hamiltonians that is uniform both in the path parameter and in the
chain length. The source obtains it from Nachtergaele's bound, whose
finite-size input is the gap of the Hamiltonian restricted to a fixed number
of sites; that gap has a uniform lower bound because the restricted
Hamiltonian depends continuously on the parameter
\cite[lines~2575--2578]{Schuch2011Phases}. This section proves the
continuity step and combines it with Knabe's nearest-neighbor criterion and
compactness. The single-block results apply to normal tensors with a common
injectivity length $p$, at every fixed interaction range $R>p$.
Quantum Wielandt supplies $p=D^4$ from normality alone. The sharper
one-site injective result has range two. The source's path also covers
normal forms with several blocks; the scope of the compact-gap conclusions
is recorded in
\cite{gap:spc11_uniform_gap_injective_scope}.

\begin{theorem}[Gap stability under perturbation of operator and kernel projection]
    \label{thm:kernel_gap_perturbation}
    \lean{ContinuousLinearMap.kernel_gap_perturbation,
        ContinuousLinearMap.eventually_norm_gap_on_orthogonal}
    \leanok
    Let $H,H_0,P,P_0$ be bounded operators on a complex normed space, and
    suppose $\delta\lVert v-P_0v\rVert\leq\lVert H_0v\rVert$ for all $v$, with
    $\delta\geq0$, $\lVert H-H_0\rVert\leq\varepsilon$ and
    $\lVert P-P_0\rVert\leq r$. Then
    $(\delta(1-r)-\varepsilon)\lVert v\rVert\leq\lVert Hv\rVert$ for every $v$
    with $Pv=0$. Consequently, if $x\mapsto H_x$ and the orthogonal projection
    onto a subspace $K_x$ are continuous at $x_0$, $K_{x_0}\subseteq\ker H_{x_0}$,
    and $\delta\lVert v\rVert\leq\lVert H_{x_0}v\rVert$ on $K_{x_0}^\perp$ with
    $\delta>0$, then every $\delta'<\delta$ satisfies
    $\delta'\lVert v\rVert\leq\lVert H_xv\rVert$ on $K_x^\perp$ for all $x$
    near $x_0$.
\end{theorem}

\begin{proof}
    \leanok
    For $Pv=0$ one has $\lVert P_0v\rVert=\lVert(P-P_0)v\rVert\leq r\lVert v\rVert$
    and $\lVert H_0v\rVert\leq\lVert Hv\rVert+\varepsilon\lVert v\rVert$;
    combining these with the hypothesis gives the first claim. For the second,
    the gap on $K_{x_0}^\perp$ and $K_{x_0}\subseteq\ker H_{x_0}$ give
    $\delta\lVert v-P_{x_0}v\rVert\leq\lVert H_{x_0}v\rVert$ for all $v$, and
    one chooses $\varepsilon=(\delta-\delta')/2$ and $r=\varepsilon/\delta$.
\end{proof}

\begin{theorem}[Open-chain gaps at every sufficient range]
    \label{thm:normal_open_gap_every_sufficient_range}
    \lean{MPSTensor.exists_openParentHamiltonianES_uniform_gap_of_isNBlkInjective_all_lengths,
        MPSTensor.exists_openParentHamiltonianES_uniform_gap_of_isNBlkInjective_of_lt,
        MPSTensor.exists_openParentHamiltonianES_uniform_gap_of_isNormal_of_le,
        MPSTensor.exists_openParentHamiltonianES_uniform_gap_of_isNormal_all_lengths}
    \leanok
    \uses{lem:open_parent_range_comparison,
        thm:open_parent_hamiltonian_kernel_ground_space}
    Let $A$ be injective at a length $p>0$, and let $R>p$. There is
    $\delta>0$ such that
    \[
        \delta\lVert v\rVert\leq\lVert H^{\mathrm{open}}_{N,R}(A)v\rVert
        \qquad(N\geq R,\ v\perp\ker H^{\mathrm{open}}_{N,R}(A)).
    \]
    In particular, for a normal tensor of bond dimension $D>0$, every
    fixed $R\geq D^4+1$ has this property. No normalization is required.
\end{theorem}

\begin{proof}
    \leanok
    \uses{lem:open_parent_range_comparison,
        thm:open_parent_hamiltonian_kernel_ground_space}
    Injectivity propagates to length $R-1$. The long-chain martingale
    estimate and range comparison give a positive gap at range $R$.
    The exact open kernel identity and finite dimensionality handle the
    finitely many shorter admissible chains. Quantum Wielandt supplies
    $p=D^4$ for a normal tensor.
\end{proof}

\begin{theorem}[Continuity of injective range projectors]
    \label{thm:continuous_injective_range_projector}
    \lean{ContinuousLinearMap.continuous_injectiveRangeProjector}
    \leanok
    Let $E$ be a finite-dimensional Hilbert space, $F$ a Hilbert space, and
    $T_x\colon E\to F$ a continuous family of injective bounded maps. Then the orthogonal projection
    $T_x(T_x^\dagger T_x)^{-1}T_x^\dagger$ onto the range of $T_x$ depends
    continuously on $x$.
\end{theorem}

\begin{proof}
    \leanok
    The Gram operator $T_x^\dagger T_x$ is invertible and continuous in $x$,
    and inversion is continuous at invertible elements of a Banach algebra.
\end{proof}

\begin{theorem}[Continuity of finite-window parent Hamiltonians]
    \label{thm:continuous_finite_window_parent_hamiltonian}
    \lean{MPSTensor.continuous_groundSpaceMapES_family,
        MPSTensor.continuous_groundSpaceES_starProjection_family,
        MPSTensor.continuous_openParentHamiltonianES_family,
        MPSTensor.continuous_parentHamiltonianES_family}
    \leanok
    \uses{def:ground_space_map_es, def:ground_space_es,
        def:open_parent_hamiltonian_es, def:parent_hamiltonian_es,
        thm:continuous_injective_range_projector}
    Let $x\mapsto A_x$ be a continuous family of tensors. For fixed $L$, the
    boundary map $X\mapsto\sum_\sigma\operatorname{tr}(A_{x,\sigma}X)\lvert\sigma\rangle$
    on $L$ sites is continuous in $x$. If every $A_x$ is $L$-block injective,
    the orthogonal projection onto the local ground space on $L$ sites is
    continuous in $x$. Hence, for $L\leq N$, the open and periodic parent
    Hamiltonians with interaction range $L$ on $N$ sites are continuous in $x$.
\end{theorem}

\begin{proof}
    \leanok
    \uses{thm:continuous_injective_range_projector,
        thm:ground_space_map_es_single,
        thm:ground_space_map_es_injective_block,
        thm:ground_space_map_es_injective_range_projector,
        lem:local_term_es_average}
    Each matrix entry of the boundary map is a trace of a fixed product of
    tensor letters. Block injectivity makes the boundary map injective, and its
    range is the local ground space, so the previous theorem applies. The
    Hamiltonians are finite sums of translates of one minus that projection.
\end{proof}

\begin{theorem}[Local stability of a finite-window gap]
    \label{thm:open_parent_gap_continuity}
    \lean{MPSTensor.eventually_openParentHamiltonianES_gap}
    \leanok
    \uses{thm:kernel_gap_perturbation,
        thm:continuous_finite_window_parent_hamiltonian}
    Let $x\mapsto A_x$ be continuous, $R\leq N$, and suppose every $A_x$ is
    $R$-block and $N$-block injective. If
    $\delta\lVert v\rVert\leq\lVert H^{\mathrm{open}}_{N,R}(A_{x_0})v\rVert$ for
    all $v\perp G_N^{\mathrm{ES}}(A_{x_0})$, with $\delta>0$, then for every
    $\delta'<\delta$ and every $x$ near $x_0$,
    $\delta'\lVert v\rVert\leq\lVert H^{\mathrm{open}}_{N,R}(A_x)v\rVert$ for
    all $v\perp G_N^{\mathrm{ES}}(A_x)$. This is the continuity step of
    \cite[lines~2575--2578]{Schuch2011Phases}.
\end{theorem}

\begin{proof}
    \leanok
    \uses{thm:kernel_gap_perturbation,
        thm:continuous_finite_window_parent_hamiltonian,
        thm:open_ground_space_le_hamiltonian_kernel}
    The open parent Hamiltonian annihilates the ground space, and both it and
    the ground-space projection are continuous in $x$.
\end{proof}

\begin{lemma}[Comparison of open-chain ranges]
    \label{lem:open_parent_range_comparison}
    \lean{MPSTensor.openParentHamiltonianES_comparison_of_local,
        MPSTensor.openParentHamiltonianES_gap_of_long_gap_at_length,
        MPSTensor.openParentHamiltonianES_gap_of_long_gap,
        MPSTensor.openParentHamiltonianES_two_gap_of_long_gap}
    \leanok
    \uses{def:open_parent_hamiltonian_es, thm:open_parent_hamiltonian_es_c1}
    Let $0<R\leq W\leq N$. If
    $\kappa\,h_W\leq H^{\mathrm{open}}_{W,R}$ for the range-$W$ parent
    interaction $h_W$ and some $\kappa>0$, then
    $\kappa H^{\mathrm{open}}_{N,W}\leq(W-R+1)H^{\mathrm{open}}_{N,R}$. For a
    tensor for which both open kernels equal $G_N^{\mathrm{ES}}$, a gap
    $\gamma\lVert v\rVert\leq\lVert H^{\mathrm{open}}_{N,W}v\rVert$ on
    $(G_N^{\mathrm{ES}})^\perp$ gives
    $\frac{\kappa\gamma}{W-R+1}\lVert v\rVert\leq\lVert H^{\mathrm{open}}_{N,R}v\rVert$
    on the same complement. The assertion may be used at one volume or
    uniformly over every $N\geq W$. For a one-site injective tensor with
    $R=2$ and $W\geq2$, the kernel hypotheses hold automatically.
\end{lemma}

\begin{proof}
    \leanok
    \uses{thm:open_parent_hamiltonian_es_c1,
        thm:open_parent_hamiltonian_kernel_ground_space,
        lem:positive_gap_transfer_loewner}
    Translate the local comparison to every window of $W$ consecutive sites and
    sum; each range-$R$ term lies in at most $W-R+1$ windows. Both open
    Hamiltonians have kernel $G_N^{\mathrm{ES}}$, so the operator inequality
    transfers the gap.
\end{proof}

\begin{theorem}[Strict Knabe windows for injective tensors]
    \label{thm:injective_strict_knabe_window}
    \lean{MPSTensor.exists_openParentHamiltonianES_two_uniform_gap_of_isInjective,
        MPSTensor.exists_strict_openParentHamiltonianES_two_window_of_isInjective}
    \leanok
    \uses{thm:primitive_open_parent_hamiltonian_gap,
        lem:parent-finite-interval-comparison,
        lem:open_parent_range_comparison}
    Let $A$ be a one-site injective tensor of nonzero bond dimension. Then the
    open-chain Hamiltonian $H^{\mathrm{open}}_{N,2}(A)$ built from the
    canonical two-site parent interaction has a positive gap above
    $G_N^{\mathrm{ES}}(A)$, uniform in all sufficiently large $N$. In
    particular, there are $m\geq2$ and $\gamma$ with $m\gamma>1$ such that
    $\gamma\lVert v\rVert\leq\lVert H^{\mathrm{open}}_{m+1,2}(A)v\rVert$ for all
    $v\perp G_{m+1}^{\mathrm{ES}}(A)$.
\end{theorem}

\begin{proof}
    \leanok
    \uses{thm:normalized_primitive_gauge,
        thm:primitive_open_parent_hamiltonian_gap,
        lem:parent-finite-interval-comparison,
        lem:open_parent_range_comparison,
        thm:open_parent_hamiltonian_kernel_ground_space}
    Gauge $A$ to a primitive tensor with a positive-definite fixed point, which
    has the same local ground spaces, and take its open-chain gap
    $\gamma_W\lVert v\rVert\leq\lVert H^{\mathrm{open}}_{N,W}(A)v\rVert$ at some
    range $W\geq2$ for all $N\geq W$. The finite-interval comparison gives
    $\kappa>0$ with $\kappa\,h_W\leq H^{\mathrm{open}}_{W,2}(A)$, and
    Lemma~\ref{lem:open_parent_range_comparison} with $R=2$ yields, for every
    $N\geq W$ and every $v\perp G_N^{\mathrm{ES}}(A)$,
    \begin{align}
        \frac{\kappa\gamma_W}{W-1}\lVert v\rVert
         & \leq\lVert H^{\mathrm{open}}_{N,2}(A)v\rVert.
        \label{eq:injective_two_site_open_gap}
    \end{align}
    Writing $\delta=\kappa\gamma_W/(W-1)$, choose $n$ with $n\delta>1$ and put
    $m=n+W$; then $m\delta>1$ and (\ref{eq:injective_two_site_open_gap}) holds
    at $N=m+1$.
\end{proof}

\begin{theorem}[Gap in a parameter neighborhood]
    \label{thm:knabe_gap_neighborhood}
    \lean{MPSTensor.eventually_parentHamiltonianES_gap_of_strict_openGap}
    \leanok
    \uses{thm:open_parent_gap_continuity,
        thm:parent_hamiltonian_gap_from_open_knabe_block}
    Let $x\mapsto A_x$ be a continuous family of one-site injective tensors of
    nonzero bond dimension, let $m\geq2$, and suppose
    $\gamma\lVert v\rVert\leq\lVert H^{\mathrm{open}}_{m+1,2}(A_{x_0})v\rVert$
    for $v\perp G_{m+1}^{\mathrm{ES}}(A_{x_0})$ with $m\gamma>1$. Then there is
    $\delta>0$ such that, for all $x$ near $x_0$ and all $N\geq2m$, the
    periodic parent Hamiltonian $H_{N,2}(A_x)$ satisfies
    $\delta\lVert v\rVert\leq\lVert H_{N,2}(A_x)v\rVert$ on the orthogonal
    complement of its kernel.
\end{theorem}

\begin{proof}
    \leanok
    \uses{thm:open_parent_gap_continuity,
        thm:parent_hamiltonian_gap_from_open_knabe_block,
        thm:open_parent_hamiltonian_kernel_ground_space}
    Choose $1/m<\gamma'<\gamma$. The strict window persists with $\gamma'$
    near $x_0$, and Knabe's nearest-neighbor criterion turns it into a periodic
    gap for every $N\geq2m$ with a constant depending only on $m$ and
    $\gamma'$.
\end{proof}

\begin{theorem}[Spectral gap along an ordered interpolation]
    \label{thm:common_kernel_spectral_gap_interpolation}
    \lean{Matrix.gap_interpolation_of_le,
        Matrix.orthogonal_quadratic_gap_of_spectrum_separated,
        Matrix.spectrum_separated_of_orthogonal_quadratic_gap,
        Matrix.spectrum_gap_interpolation_of_le,
        Matrix.spectrum_gap_interpolation_of_le_of_nontrivial_kernel}
    \leanok
    If every eigenvalue $z$ of a Hermitian matrix $M$ has $\operatorname{Re}z=0$
    or $\operatorname{Re}z\geq\delta$, then
    $\delta\lVert x\rVert^2\leq\operatorname{Re}\langle x,Mx\rangle$ for all
    $x\perp\ker M$. Conversely, if $M\geq0$ satisfies this quadratic-form bound,
    then every eigenvalue $z$ of $M$ has $\operatorname{Re}z\geq0$, and
    $\operatorname{Re}z=0$ or $\operatorname{Re}z\geq\delta$.
    Let $0\leq A\leq B$ be matrices with $\ker A\subseteq\ker B$ and such a
    gap $\delta$ for $A$, and let $t\geq0$. Then $C_t=A+t(B-A)$ is positive
    semidefinite, $\ker C_t=\ker A$, and every eigenvalue $z$ of $C_t$
    satisfies $\operatorname{Re}z=0$ or $\operatorname{Re}z\geq\delta$. If
    $\ker A\neq0$, then $0$ is an eigenvalue of every $C_t$.
\end{theorem}

\begin{proof}
    \leanok
    Since $t(B-A)\geq0$, a vector annihilated by $C_t$ has
    $\langle x,Ax\rangle=0$, hence $Ax=0$, and conversely $Ax=0$ forces $Bx=0$.
    On $(\ker A)^\perp$ one has
    $\langle x,C_tx\rangle\geq\langle x,Ax\rangle\geq\delta\lVert x\rVert^2$. The
    equivalence with the spectral statement follows by expanding $x$ in an
    orthonormal eigenbasis of the Hermitian matrix.
\end{proof}

\begin{lemma}[Uniform bounds from local bounds on a compact set]
    \label{lem:compact_uniform_pos_nat_bounds}
    \lean{IsCompact.exists_uniform_pos_nat_bounds}
    \leanok
    Let $S$ be a compact subset of a topological space and $P(\delta,N,x)$ a
    property of a real $\delta$, a natural number $N$ and a point $x$ that is
    preserved by decreasing $\delta$ and by increasing $N$. If every $x\in S$
    has $\delta>0$ and $N$ such that $P(\delta,N,y)$ holds for all $y$ near $x$,
    then one $\delta>0$ and one $N$ satisfy $P(\delta,N,x)$ for all $x\in S$.
\end{lemma}

\begin{proof}
    \leanok
    Take a finite subcover of the neighborhoods, the minimum of the finitely
    many constants and the maximum of the finitely many thresholds.
\end{proof}

\begin{theorem}[Uniform finite-dimensional kernel gap]
    \label{thm:compact_kernel_gap}
    \lean{ContinuousLinearMap.exists_uniform_norm_gap_of_compact}
    \leanok
    \uses{thm:kernel_gap_perturbation, lem:compact_uniform_pos_nat_bounds}
    Let $E$ be a finite-dimensional complex Hilbert space, let
    $x\mapsto H_x\colon E\to E$ be continuous, and suppose the orthogonal
    projections onto $\ker H_x$ are continuous in $x$. For every compact
    parameter set $S$, there is $\delta>0$ such that
    $\delta\lVert v\rVert\leq\lVert H_xv\rVert$ for every $x\in S$ and
    $v\perp\ker H_x$. No positivity or self-adjointness is required.
\end{theorem}

\begin{proof}
    \leanok
    \uses{thm:kernel_gap_perturbation, lem:compact_uniform_pos_nat_bounds}
    The restriction of $H_x$ to $(\ker H_x)^\perp$ is injective, so finite
    dimensionality gives a positive lower norm bound at each parameter.
    Half of that bound persists in a neighborhood by continuity of the
    operator and its kernel projection. A finite subcover of $S$ gives one
    positive bound for all parameters.
\end{proof}

\begin{theorem}[Continuity of short periodic ground lines]
    \label{thm:periodic_short_gap_continuity}
    \lean{MPSTensor.exists_uniform_parentHamiltonianES_two_gap_fixed_volume}
    \leanok
    \uses{thm:compact_kernel_gap, thm:continuous_injective_range_projector,
        thm:continuous_finite_window_parent_hamiltonian}
    Let $x\mapsto A_x$ be a continuous family of one-site injective tensors of
    nonzero bond dimension, let $S$ be a compact parameter set, and let
    $N\geq2$. Then there is $\delta>0$ such that
    $\delta\lVert v\rVert\leq\lVert H_{N,2}(A_x)v\rVert$ for all $x\in S$ and
    all $v$ orthogonal to the kernel of $H_{N,2}(A_x)$.
\end{theorem}

\begin{proof}
    \leanok
    \uses{thm:compact_kernel_gap, thm:continuous_injective_range_projector,
        thm:continuous_finite_window_parent_hamiltonian}
    The kernel of $H_{N,2}(A_x)$ is the line spanned by the nonzero periodic
    matrix product vector of $A_x$, whose projection is continuous in $x$. The
    gap at each point therefore persists in a neighborhood, and a finite
    subcover of $S$ gives one constant.
\end{proof}

\begin{theorem}[Uniform parent gap on compact parameter sets]
    \label{thm:compact_parent_gap}
    \lean{MPSTensor.exists_uniform_parentHamiltonianES_gap_of_compact_isInjective,
        MPSTensor.exists_uniform_parentHamiltonianES_gap_of_compact_isInjective_all_lengths}
    \leanok
    \uses{thm:knabe_gap_neighborhood, thm:injective_strict_knabe_window,
        thm:periodic_short_gap_continuity, lem:compact_uniform_pos_nat_bounds}
    Let $x\mapsto A_x$ be a continuous family of one-site injective tensors of
    nonzero bond dimension, and let $S$ be a compact parameter set. Then there
    is $\delta>0$ such that, for every $x\in S$, every $N\geq2$ and every
    $v\perp\ker H_{N,2}(A_x)$,
    \begin{align}
        \delta\lVert v\rVert & \leq\lVert H_{N,2}(A_x)v\rVert.
        \label{eq:compact_parent_gap}
    \end{align}
    This is the uniform-gap conclusion of
    \cite[Appendix~A, lines~2575--2580]{Schuch2011Phases} for one-site
    injective tensors and the two-site parent interaction; the source also
    treats non-injective normal forms
    \cite{gap:spc11_uniform_gap_injective_scope}.
\end{theorem}

\begin{proof}
    \leanok
    \uses{thm:knabe_gap_neighborhood, thm:injective_strict_knabe_window,
        thm:periodic_short_gap_continuity, lem:compact_uniform_pos_nat_bounds}
    Every point of $S$ has a strict Knabe window, hence a neighborhood with a
    common gap bound for all $N$ beyond a threshold. A finite subcover of $S$
    gives one bound and one threshold. The finitely many shorter rings are
    handled by the previous theorem, and the minimum of these constants works
    for every $N\geq2$.
\end{proof}

\begin{theorem}[Compact normal families at a fixed interaction range]
    \label{thm:compact_normal_parent_gap}
    \lean{MPSTensor.exists_openParentHamiltonianES_uniform_gap_of_isNBlkInjective,
        MPSTensor.exists_strict_openParentHamiltonianES_window_of_isNBlkInjective,
        MPSTensor.eventually_parentHamiltonianES_gap_of_strict_openGap_of_isNBlkInjective,
        MPSTensor.exists_uniform_parentHamiltonianES_gap_of_compact_isNBlkInjective,
        MPSTensor.exists_uniform_parentHamiltonianES_gap_of_compact_isNBlkInjective_all_lengths,
        MPSTensor.exists_uniform_parentHamiltonianES_gap_fixed_volume_of_isNBlkInjective,
        MPSTensor.isNBlkInjective_pow_four_of_isNormal,
        MPSTensor.periodicMpvLineMap_injective_of_isNBlkInjective,
        MPSTensor.ker_parentHamiltonianES_eq_range_periodicMpvLineMap_of_isNBlkInjective,
        MPSTensor.continuous_parentHamiltonianES_kernelProjection_family_of_isNBlkInjective,
        MPSTensor.exists_uniform_parentHamiltonianES_gap_of_compact_isNormal_all_lengths,
        MPSTensor.exists_uniform_parentHamiltonianES_gap_of_compact_isNBlkInjective_at_range,
        MPSTensor.exists_uniform_parentHamiltonianES_gap_of_compact_isNormal_at_range}
    \leanok
    \uses{thm:compact_kernel_gap, lem:open_parent_range_comparison,
        thm:parent_hamiltonian_gap_from_open_knabe_block}
    Let $x\mapsto A_x$ be continuous with nonzero fixed bond dimension,
    and let $S$ be compact. Suppose every $A_x$ is injective after the same
    positive length $p$. For every fixed interaction range $R>p$, there is
    $\delta>0$ such that
    \[
        \delta\lVert v\rVert\leq\lVert H_{N,R}(A_x)v\rVert
        \qquad(x\in S,\ N\geq R,\ v\perp\ker H_{N,R}(A_x)).
    \]
    If each $A_x$ is merely normal, the same assertion holds at the
    ranges $R\geq D^4+1$. No gap estimate or normalized gauge is
    supplied as a hypothesis.
\end{theorem}

\begin{proof}
    \leanok
    \uses{thm:normalized_primitive_gauge,
        thm:primitive_open_parent_hamiltonian_gap,
        lem:open_parent_range_comparison,
        thm:open_parent_gap_continuity,
        thm:parent_hamiltonian_gap_from_open_knabe_block,
        lem:compact_uniform_pos_nat_bounds, thm:compact_kernel_gap,
        thm:chain_ground_space_eq_mpv_blk,
        thm:exact_pow_four_block_injectivity_normal_left_canonical}
    A pointwise primitive gauge supplies an open-chain gap at a long range.
    The finite-interval comparison transfers it to range $p+1$. Choosing
    a sufficiently long window makes the finite-range Knabe inequality
    strict. This inequality persists near each parameter, and a finite
    subcover gives a uniform gap for all sufficiently long rings.
    The shorter rings have the line generated by their nonzero periodic
    matrix product vector as kernel. Its projection is continuous, so
    compactness supplies positive bounds for those finitely many lengths.
    For a normal tensor, quantum Wielandt gives injectivity at $D^4$ after
    a primitive gauge; gauge invariance gives the same length for the
    original tensor. For a larger prescribed range $R>p$, injectivity
    propagates to length $R-1$, and the same argument applies.
\end{proof}


\begin{theorem}[Continuity for several normal blocks]
    \label{thm:block_ground_space_continuity}
    \lean{MPSTensor.blockGroundSpaceMapES,
        MPSTensor.blockGroundSpaceMapES_apply,
        MPSTensor.continuous_blockGroundSpaceMapES_family,
        MPSTensor.blockGroundSpaceMapES_injective_of_wordTupleSpanTop,
        MPSTensor.range_blockGroundSpaceMapES,
        MPSTensor.continuous_iSup_groundSpaceES_starProjection_family,
        MPSTensor.continuous_groundSpaceES_toTensorFromBlocks_starProjection_family,
        MPSTensor.continuous_norm_groundSpaceES_overlap_family,
        MPSTensor.blockPeriodicMpvMapES,
        MPSTensor.blockPeriodicMpvMapES_apply,
        MPSTensor.blockPeriodicMpvMapES_eq_blockGroundSpaceMapES,
        MPSTensor.continuous_blockPeriodicMpvMapES_family,
        MPSTensor.ker_parentHamiltonianES_toTensorFromBlocks_eq_range_blockPeriodicMpvMapES,
        MPSTensor.blockPeriodicMpvMapES_injective_of_wordTupleSpanTop,
        MPSTensor.range_blockPeriodicMpvMapES,
        MPSTensor.continuous_parentHamiltonianES_toTensorFromBlocks_kernelProjection_family,
        MPSTensor.exists_uniform_parentHamiltonianES_toTensorFromBlocks_gap_fixed_volume_of_wordTupleSpanTop}
    \leanok
    \uses{thm:compact_kernel_gap}
    Let $x\mapsto(A_x^j)_j$ be a continuous family of blocks. At every
    length $L$ at which the length-$L$ word tuples span the product of the
    block matrix algebras for every parameter, the projection onto the sum
    of the local block spaces on $L$ sites, and the projection onto the
    local space of the weighted block sum, are continuous in the parameter.
    Suppose moreover that the physical dimension and all bond dimensions
    are positive, and that the word tuples span at a length $S>0$ for every
    parameter. For the weighted block sum, the periodic kernel projection
    is continuous whenever $S+1\leq R\leq N$. Nonzero block weights may
    vary arbitrarily. Thus each such fixed volume has a positive gap
    uniform on any compact parameter set. The overlaps of two local block spaces are also
    continuous at any common injectivity length.
\end{theorem}

\begin{proof}
    \leanok
    \uses{thm:compact_kernel_gap, thm:continuous_injective_range_projector,
        thm:block_word_span_propagation}
    The joint boundary map has fixed domain, is injective, and has range
    equal to the sum of the local block spaces. Its range projection
    therefore depends continuously on the parameter. Restricting the
    boundary matrices to scalar multiples of the identity gives the
    injective map onto the periodic ground space. The kernel identity
    then gives continuity of the periodic kernel projection, after the
    span at $S$ is propagated to every length at least $S$ using the
    positive dimensions. Compactness
    yields a common positive finite-volume gap.
\end{proof}

\begin{theorem}[Compact families of several normal blocks]
    \label{thm:compact_block_parent_gap}
    \lean{MPSTensor.eventually_parentHamiltonianES_toTensorFromBlocks_gap_of_strict_openGap,
        MPSTensor.exists_uniform_parentHamiltonianES_toTensorFromBlocks_gap_of_compact,
        MPSTensor.exists_uniform_parentHamiltonianES_toTensorFromBlocks_gap_of_compact_all_lengths,
        MPSTensor.exists_uniform_block_parentHamiltonianES_gap_of_compact_isNormalTensor}
    \leanok
    \uses{thm:block_ground_space_continuity,
        thm:cpsv_bnt_sharp_span_ambient_cap, thm:block_word_span_propagation,
        thm:block_strict_knabe_window_simultaneous_injectivity,
        thm:parent_hamiltonian_gap_from_open_knabe_block}
    Under the preceding simultaneous spanning hypothesis, fix
    $R\geq S+1$. On a compact parameter set there is $\delta>0$ such
    that the canonical periodic parent Hamiltonian of the weighted block
    sum satisfies
    \[
        \delta\lVert v\rVert\leq\lVert H_{N,R}(x)v\rVert
        \qquad(N\geq R,\ v\perp\ker H_{N,R}(x)).
    \]
    In particular, for pairwise inequivalent normalized normal blocks of
    dimensions $D_j>0$, one may choose any
    $R\geq3\max(\sum_jD_j,1)^5$. The simultaneous spanning length
    is then derived from the dimensions.
    This includes the two-site interaction of the initially blocked path
    in \cite[Appendix~A]{Schuch2011Phases}; see
    \cite{gap:spc11_uniform_gap_injective_scope}.
\end{theorem}

\begin{proof}
    \leanok
    \uses{thm:block_ground_space_continuity,
        thm:cpsv_bnt_sharp_span_ambient_cap, thm:block_word_span_propagation,
        thm:block_strict_knabe_window_simultaneous_injectivity,
        thm:parent_hamiltonian_gap_from_open_knabe_block}
    The open-chain multiblock gap gives a strict finite-range Knabe window
    at each parameter. Continuity preserves the strict inequality in a
    neighborhood. A finite subcover gives a common gap for all sufficiently
    long rings. The preceding fixed-volume result handles the finitely
    many remaining lengths, including $N=R$. For normalized normal blocks,
    the simultaneous word-span bound supplies a common length
    $S=3\max(\sum_jD_j,1)^5-1$.
\end{proof}

\begin{theorem}[Continuous positive parent interactions]
    \label{thm:compact_positive_parent_interaction_gap}
    \lean{MPSTensor.exists_uniform_pos_parentInteractionES_le_of_compact,
        MPSTensor.exists_uniform_periodicInteractionHamiltonianES_gap_of_compact_canonical_gap,
        MPSTensor.exists_uniform_parentInteraction_gap_of_compact_isNormal,
        MPSTensor.exists_uniform_block_parentInteraction_gap_of_compact,
        MPSTensor.exists_uniform_parentInteraction_gap_of_compact_isNormal_at_range}
    \leanok
    \uses{thm:compact_kernel_gap, thm:compact_normal_parent_gap,
        thm:compact_block_parent_gap}
    Suppose the local ground-space projection and a positive parent
    interaction $h_x$ depend continuously on $x$, with the same local
    kernel. On a compact parameter set there is $\kappa>0$ such that
    $h_x\geq\kappa P_{G_R(x)^\perp}$. Consequently every canonical
    periodic gap uniform in the parameter and chain length transfers to
    the chosen interactions. This applies, in particular, to compact
    normal families at every fixed range $R\geq D^4+1$.
\end{theorem}

\begin{proof}
    \leanok
    \uses{thm:compact_kernel_gap}
    Compactness gives a common positive local gap away from the kernel.
    Positivity turns it into the stated order bound. Translation and
    summation preserve this comparison. A pointwise upper comparison
    identifies the periodic kernels, and the lower comparison transfers
    the canonical gap with factor $\kappa$.
\end{proof}

